\subsection{Integral bar integration with a nonclosed degree-zero factor}

\paragraph{A-S09-13 (Proved; the factor in A-L05 need not be closed).}
Let $k$ be any commutative ring, let $r\geq1$, and let
$(R,J,\gamma)$ be a strict PD cochain $k$-algebra.
Let $M$ be a free $k$-module, let
$\alpha:M\to J^{2r-1}$ be linear, and let $z\in J^0$.
No condition is imposed on $dz$.
There is a natural integral strict PD cochain map
\[
       \mathcal I_{z,\alpha}:U_r(M)\longrightarrow R
 \tag{S09.31}
\]
whose original disk has images
$x(a)\mapsto z\alpha(a)$ and
$y(a)\mapsto d(z\alpha(a))$.

Here is a formula for the whole map.  Put $B=\Gamma_k(M)$,
give its generators form degree $2r$, and let
$\phi_0:B\to R$ be the PD map determined by
$\phi_0(a)=d\alpha(a)$.  Let $\delta_0$ be its special
derivation with $\delta_0(a)=\alpha(a)$.
For an arity-homogeneous word $b_0[b_1|\cdots|b_n]$, put
\[
 a_i=\operatorname{arity}(b_i),\quad
 m=\sum_{i=0}^na_i,\quad
 A_i=a_i+\cdots+a_n,\quad D=\prod_{i=1}^n A_i.
\]
The conventions are $a_0\geq0$, $a_i>0$ for $i>0$,
and $D=1$ when $n=0$.  The value of (S09.31) is
\[
\begin{split}
 &\frac{m!}{D}\,\gamma_m(z)\,
       \phi_0(b_0)\delta_0(b_1)\cdots\delta_0(b_n)\\
 &\qquad+
 \frac{(m-1)!}{D}\,\gamma_{m-1}(z)\,dz\,
       \delta_0(b_0)\delta_0(b_1)\cdots\delta_0(b_n).
\end{split}
\tag{S09.32}
\]
The second line is included only when $a_0>0$; when
$a_0=0$ it is defined to be zero, including the scalar
word $m=n=0$.  Every displayed coefficient that is used
is an integer.  The map is natural in linear parameter
maps and strict PD differential target maps.  It is
simplicial when $M,R,z,\alpha$ are simplicial data with
the stated compatibility.  If $z$ has weight $a$ and
$\alpha(M)$ has weight $b$, the map preserves weights
when the original disk is placed in weight $a+b$.

\begin{proof}
The positive suffix sums are strictly decreasing
integers at most $m$, so their product divides $m!$.
If $a_0>0$, they are all at most $m-1$, so their
product also divides $(m-1)!$.  This proves integrality
without division in $k$.

We prove the required polynomial identities in the free
integral strict PD differential algebra
\[
 \Gamma_{\mathbb Z}(z,\{y_a\})\otimes
       \Lambda_{\mathbb Z}(u,\{x_a\}),
 \qquad dz=u,\quad dx_a=y_a,\quad du=dy_a=0,
 \tag{S09.33}
\]
where $|z|=0$, $|u|=1$, $|x_a|=2r-1$, and $|y_a|=2r$.
The indices range over a basis of the finitely many
parameters occurring in an expression.  Its underlying
group is torsion free.  All identities proved integrally
there specialize to $R$, including when $k$ has torsion.

Let $\phi'$ and $\delta'$ be the PD map and special
derivation determined by
\[
 \phi'(a)=u x_a+z y_a=d(zx_a),\qquad
 \delta'(a)=z x_a.
\]
For a homogeneous $b\in B$ of arity $a>0$, their exact
integral identities are
\[
 \begin{split}
 \phi'(b)&=z^a\phi_0(b)+z^{a-1}u\delta_0(b),\\
 \delta'(b)&=z^a\delta_0(b).
 \end{split}
 \tag{S09.34}
\]
For a single divided power the first formula follows
from the PD sum formula: the higher divided powers of
$u x_a$, a product of two odd elements, vanish.  The
formula for a product follows by expansion, since
$u^2=0$.  For the second formula, expand the special
derivation of a divided-power monomial.  The terms
without $u$ give the stated expression.  A term with
two distinct selected odd variables has a partner
with those selections exchanged; moving $u$ to the
front gives opposite signs.  These terms cancel.
A term selecting the same odd variable twice is zero
by its alternating square.  This proves (S09.34),
including arbitrary divided-power multiplicities.

In the rationalization of (S09.33), the usual bar
integration formula for the disk $a\mapsto z x_a$ is
\[
 \frac{1}{D}\,\phi'(b_0)
          \delta'(b_1)\cdots\delta'(b_n).
 \tag{S09.35}
\]
For completeness, it is a cochain algebra map because
$d\delta'(b)=a\phi'(b)$ in arity $a$, differentiation
of the product in (S09.35) gives the coefficients
$(-1)^{j-1}a_j/D$, and the two adjacent bar mergers
give the identical coefficients
$(-1)^{j-1}(A_j-A_{j+1})/D$.
The endpoint merger has the same value, with
$A_{n+1}=0$.  Multiplicativity is the labelled-shuffle
identity
\[
 \sum_{\sigma\in\operatorname{Sh}(n,\ell)}
       \frac{1}{D(\sigma(a,b))}
              =\frac{1}{D(a)D(b)}.
\]
Partitioning the shuffles by their first entry proves
this by induction: the two contributions have
denominators beginning with the common sum
$\sum a_i+\sum b_i$, and their numerators add to
that sum.  The shuffle sign is canceled by the sign
of reordering the odd $\delta'$ factors.  These are
exactly the bar differential and shuffle multiplication
of $U_r(M)$.  Over $\mathbb Q$, the algebra map also
preserves the even divided powers $v^j/j!$.

Substituting (S09.34) into (S09.35) gives its two terms
$z^m/D$ and $z^{m-1}/D$ with the factors displayed
in (S09.32).  Replacing those powers of $z$ by their
integral PD expressions gives exactly (S09.32).
The second term is zero when $b_0$ is scalar, as stipulated.
The integrality already proved and injectivity of
(S09.33) into its rationalization now prove the chain,
product, and PD identities integrally.  Specialization
proves them over the arbitrary ring $k$.

The scalar word maps to one.  Every positive-arity
word maps into $J$: the first term has its factor
$\gamma_m(z)$, and the second has $\gamma_{m-1}(z)$
or, when $m=1$, the factor $dz\in J$.
The original bar generator $1[a]$ maps to $z\alpha(a)$;
the coefficient generator $a[]$ maps to
$z\,d\alpha(a)+dz\,\alpha(a)$.
Both terms in (S09.32) have form degree $2rm-n$ and
weight $m(a+b)$ under the stated weight assignment.
Finally $\phi_0,\delta_0$, the PD operations, and $d$
commute with the stated parameter and target maps,
while the suffix coefficients depend only on arity.
This proves naturality and its simplicial consequence.
\end{proof}

For one parameter, the primitive already used in A-L04
therefore works without the closedness assumption:
\[
 d\bigl((n-1)!\gamma_n(z)x\gamma_{n-1}(dx)\bigr)
                    =\gamma_n\bigl(d(zx)\bigr).
\]
The additional derivative of $z$ is precisely the
single $dz\,x$ term in the divided power on the right.
This construction supplies full boundary maps for the
stated factorizations.  No factorization of the general
cotangent classes or of an arbitrary prescribed
boundary map is part of the assertion.
